\section{Experiment Plan}
\subsection{Applications Under Test} \label{sec:apps}
The two microservice applications used to test the accuracy of the simulators are \textit{Sock Shop} \cite{weaveworks_sock_2025} (see \autoref{fig:sock}) and \textit{Train Ticket} \cite{codewisdom_train_2025} (see \autoref{fig:train}). An overview of properties that can be used to describe these applications can be found in \autoref{tab:microservices}. Both simulate e-commerce application for online stores consisting of services that communicate with each other via REST API over HTTP. While we can see from the number of services and different programming languages used that \textit{Train Ticket} is more complex. \textit{Sock Shop}, on the other hand, is far more popular and regularly used in research, including recent research into topics related to the energy consumption of microservices \cite{floroiu_anomaly_2024, rao_energy-aware_2025}.  

\begin{table}[ht]
    \centering
    \caption{Description of the microservices applications used for testing.\tablefootnote{Retrieved 20 June 2025. \label{retrievaldate}}}
    \begin{tabular}{l|ll}
    \textbf{Property} & \textit{\textbf{Sock Shop}} & \textit{\textbf{Train Ticket}} \\ \hline
    Created & 2016 & 2019 \\ \arrayrulecolor{lightgray}\hline
    Last updated & 2023 & 2022 \\ \arrayrulecolor{lightgray}\hline
    Number of services & 13 & 59 \\ \arrayrulecolor{lightgray}\hline
    \begin{tabular}[c]{@{}l@{}}Programming\\languages\end{tabular} & Java, JavaScript, Go & \begin{tabular}[c]{@{}l@{}}Java, JavaScript,\\Python, Go\end{tabular} \\ \arrayrulecolor{lightgray}\hline
    Types of databases & Mongo, MySQL & Mongo, MySQL \\ \arrayrulecolor{lightgray}\hline
    \begin{tabular}[c]{@{}l@{}}Communication\\protocol\end{tabular} & HTTP & HTTP \\ \arrayrulecolor{lightgray}\hline
    Commits & 1705 & 323 \\ \arrayrulecolor{lightgray}\hline
    Contributors & 55 & 26 \\ \arrayrulecolor{lightgray}\hline
    Forks & 2875 & 272 \\ \arrayrulecolor{lightgray}\hline
    Stars & 3702 & 788
    \end{tabular}
    \label{tab:microservices}
\end{table}

\begin{figure}[ht]
    \centering
    \includegraphics[width=0.5\linewidth]{vu-cs-research-thesis/figures/Microservice systems-Sock Shop.png}
    \caption{System architecture of Sock Shop.}
    % \Description{System architecture of Sock Shop.}
    \label{fig:sock}
\end{figure}

\begin{figure}[ht]
    \centering
    \includegraphics[width=0.9\linewidth]{vu-cs-research-thesis/figures/Microservice systems-Train Ticket.png}
    \caption{System architecture of Train Ticket.}
    % \Description{System architecture of Train Ticket.}
    \label{fig:train}
\end{figure}


\subsection{Simulators Under Test} \label{sec:simulators}

\subsubsection{Simulator Selection}
In order to select the simulators analysed in this study, a systematic selection is carried out (see \autoref{fig:selection}). First, 23 papers are collected by querying Google Scholar\footnote{\url{https://scholar.google.com}} for \textit{Cloud computing simulators review} with search period 2019-present, \textit{Cloud computing simulator}, and \textit{“cloud computing simulators”} with search period 2020-present. The results are ordered by relevance and the top results are collected\footnote{This was done on 4 February 2025.}. Furthermore, the forward snowballing method \cite{wohlin_experimentation_2012} is used by collecting papers that cite \cite{mansouri_cloud_2020}\footnote{\url{https://scholar.google.com/scholar?cites=8172644899227091971}}, as this is the most influential relevant paper obtained through the previous queries having been cited more than 100 times. This results in a total of 23 papers, 13 of which are determined to be relevant as discussing multiple cloud simulators and their properties (see \autoref{tab:papersAppendix}). From the 13 papers, all simulators mentioned are extracted as well as 10 simulators that are suggested by a research practitioner in the field of microservices. This leads to a total of 67 simulators for which the date of the first publication presenting the simulator and the date of the most recent update to the simulator are retrieved (see \autoref{tab:allSimulatorsAppendix}). 

\begin{figure}[htbp]
    \centering
    \includegraphics[width=0.5\linewidth]{vu-cs-research-thesis/figures/selection.png}
    \caption{Simulator selection process.}
    % \Description{Simulator selection process.}
    \label{fig:selection}
\end{figure}

Furthermore, each simulator is assessed based on whether its code is publicly available and whether the simulator has energy usage simulation capabilities. Only simulators for which the code is publicly available are considered both to increase the replicability of this study as well as to be able to inspect their code. 37 simulators fulfil both these requirements. From the 37 simulators that fulfil both these requirements, the 6 simulators that have been most recently updated are selected: \textit{CloudSim}, \textit{NetworkCloudSim}, \textit{OpenDC}, \textit{SimGrid}, and \textit{DISSECT-CF}. \autoref{tab:simulators-overview} gives an overview of their properties. While \textit{CloudSim Plus} was also recently updated, it was not included in this study as it runs on the \textit{CloudSim} engine, which is already included in the study through both \textit{CloudSim} and \textit{NetworkCloudSim}. Out of the five selected simulators, \textit{CloudSim} is by far the most popular simulator with the largest number of forks, stars and citations, while \textit{OpenDC} is the newest simulator, having first been published in 2021. In remainder of this section, a general description of each simulator is given.

\begin{table*}[ht]
\centering
\caption{The selected cloud simulators\protect\footnotemark[2].}
\label{tab:simulators-overview}
\begin{tabular}{l|lllll}
\textbf{Property}                                                          & \textbf{CloudSim} & \textbf{DISSECT-CF} & \textbf{OpenDC} & \textbf{NetworkCloudSim} & \textbf{SimGrid}                                                   \\ 
\hline
Created                                                                   & 2009              & 2015                & 2021            & 2011                     & 2001                                                               \\ \arrayrulecolor{lightgray}\hline
Last updated                                                              & 2025              & 2023                & 2025            & 2025                     & 2025                                                               \\ \arrayrulecolor{lightgray}\hline
Based on                                                                  & \textit{SimJava}, \textit{GridSim}           &  -                   &  -               & \textit{CloudSim}                 & -                                                                \\ \arrayrulecolor{lightgray}\hline
Commits                                                                   & 465 \cite{cloudslab_cloudsim_2025}               & 466 \cite{kecskemeti_dissect-cf_2023}                & 2267 \cite{mastenbroek_opendc_2021}           & 465 \cite{cloudslab_cloudsim_2025}                    & 37573  \cite{simgrid_simgrid_2025}                                                            \\ \arrayrulecolor{lightgray}\hline
Contributors                                                              & 15                & 4                   & 17              & 15                       & 70                                                                   \\ \arrayrulecolor{lightgray}\hline
Forks                                                                     & 532               & 44                  & 62              & 532                      &  179                                                                  \\ \arrayrulecolor{lightgray}\hline
Stars                                                                     & 912               & 6                   & 86              & 912                      & 94                                                                  \\ \arrayrulecolor{lightgray}\hline
\begin{tabular}[c]{@{}l@{}}Programming\\languages simulator\end{tabular}  & Java              & Java                & Kotlin          & Java                     & \begin{tabular}[c]{@{}l@{}}C++, C, \\Python, or Java\end{tabular}  \\ \arrayrulecolor{lightgray}\hline
\begin{tabular}[c]{@{}l@{}}Programming\\languages simulation\end{tabular} & Java              & Java                & -               & Java                     & \begin{tabular}[c]{@{}l@{}}C++, C,\\Python, or Java\end{tabular}   \\ \arrayrulecolor{lightgray}\hline
\begin{tabular}[c]{@{}l@{}}Citations first paper\footnotemark[3] \end{tabular}         & 1724              & 125                 & 37              & 446                      & 695                                                      
\end{tabular}
\end{table*}

\setcounter{footnote}{3}

\subsubsection{CloudSim}
\textit{CloudSim} \cite{cloudsim_buyya_modeling_2009} is a simulation toolkit for simulating cloud environments, built on the \textit{SimJava} \cite{howell_simjava_1998} discrete event simulation engine and the \textit{GridSim} \cite{buyya_gridsim_2002} Grid simulation toolkit. It was first published in 2009 and is still being developed by the CLOUDS Laboratory at the University of Melbourne. It is written in Java and implements \verb|Brokers| and \verb|Datacenters|, with the latter being assigned stand-alone objects such as \verb|Hosts|, \verb|VMs| and processing units called \verb|Cloudlets|. For the allocation and scheduling of tasks, provided allocation and scheduling policies can be used, or they can be custom defined by the user. Workloads can be loaded from files in \textit{Standard Workload Format} \cite{chapin_swf_1999} (swf) or defined by the user in code.

The capability to simulate energy usage within CloudSim was added in 2012 \cite{cloudsimPower_beloglazov_optimal_2012}, with dedicated power objects such as \\ \verb|PowerDatacenter| and \verb|PowerHost|. For the calculation of the simulated energy usage, only the hosts' CPU usage is taken into account. Four power models (cubic, linear, square root and square) are provided, and the user can define a custom power model.

\subsubsection{DISSECT-CF}
\textit{DISSECT-CF} \cite{kecskemeti_dissect-cf_2015} is an IaaS simulation framework, first published in 2015. It was developed with a focus on modelling energy usage and evaluating different scheduling strategies. It is written in Java and contains five subsystems of decreasing level of abstraction: an event system, unified resource sharing, energy modelling, infrastructure simulation and infrastructure management. Workloads can be loaded from files in gwf, swf, srtg or one2 format, or defined in code. For the simulation of energy usage, \textit{DISSECT-CF} decouples energy consumption from resource utilization. Instead, it uses resource utilisation counters that are updated depending on the power state of a resource. For each individual power state, a different power consumption model is defined. Predefined power states for a physical machine are \textit{off}, \textit{switching on}, \textit{running} and \textit{switching off}. \textit{DISSECT-CF} provides three predefined power models, constant, linear and NoIdle, a variation on the linear model, as well as the possibility for the user to define their own custom power model.

\subsubsection{NetworkCloudSim}
\textit{NetworkCloudSim} \cite{garg_networkcloudsim_2011} is an extension to the \textit{CloudSim} framework that enables the simulation of network components and message passing, first published in 2011. It is also implemented in Java with dedicated network objects, such as \verb|NetworkDatacenter| and \verb|NetworkHost|. In their original implementation, these objects are incompatible with the dedicated power objects, but their capabilities do not interfere with each other. Because of this, we combined the power and network elements where necessary and created for example a \verb|PowerNetworkHost| in order to still be able to evaluate \textit{NetworkCloudSim}. As we only combined existing objects and did not add any energy usage simulation capabilities, this means that also for \textit{NetworkCloudSim} only CPU usage is taken into account and the same four power models are available as for \textit{CloudSim}, next to the possibility for the user to define a custom power models. As is the case for \textit{CloudSim}, workloads can be given in swf format or defined by the user. Communication is defined as actions between \verb|Cloudlets| and have to be defined in code.

\footnotetext[2]{See footnote \ref{retrievaldate}.}
\footnotetext[3]{Obtained using the Google Scholar "cited by" function.}

\subsubsection{OpenDC}
\textit{OpenDC} \cite{mastenbroek_opendc_2021} is a platform for datacenter simulation, first published in 2021. It is implemented in Kotlin, but experiments on the platform and their topologies are supplied as JSON files. This leads to the platform having a high level of abstraction and a low level of customizability. Simulated workloads are given as parquet\footnote{\url{https://parquet.apache.org/}} files. To calculate the simulated power usage, only CPU usage is taken into account and \textit{OpenDC} provides five power models: constant, which is the default, cubic, linear, square root and square. It is not possible for users to define their own power models. OpenDC does provide the option to include carbon traces, which are then used to calculate the carbon intensity alongside the energy usage.

\subsubsection{SimGrid}
\textit{SimGrid} \cite{casanova_simgrid_2025} is a parallel and distributed computing (PDC) simulator, first published in 2001\cite{casanova_simgrid_2001}. It is available in C++, C, Python and Java, we used the version written in C++ for this study as it is the most extensive. Workloads have to be defined in code. The energy usage simulation capabilities of SimGrid are defined in the \textit{Host Energy} plugin \cite{simgridPower_heinrich_predicting_2017} and \textit{Link Energy} plugin. Next to these plugins, there are also \textit{WiFi Energy}, \textit{Battery}, \textit{Solar Panel} and \textit{Chiller} plugins, but these will not be studied here. Energy consumption is modelled linearly, based on off, idle, maximum utilization, and an \verb|Epsilon| value for a near-idle state. Energy consumption for a certain CPU usage is calculated through linear extrapolation between the \verb|Epsilon| and maximum values.

\subsection{Experiment Variables}
\subsection{Independent Variables} \label{sec:independentVars}
For both RQs, the subjects are the cloud simulators (see \autoref{sec:simulators}) tested. For these simulators, various different simulation configurations are tested (see \autoref{tab:config-vars}).

For \textit{CloudSim} and \textit{NetworkCloudSim}, workloads are defined in code with one \verb|Cloudlet| per service. For \textit{SimGrid}, an actor is created for each service which is then set to execute computations according to the CPU utilization from the workload's trace file.
For each simulator, all available predefined power models are tested. Custom power models are not implemented and then tested as these depend too much on the users' implementation and are therefore not representative of the simulator, which is the main object of study.

\begin{table*}[ht]
    \centering
    \caption{Simulation configuration variables tested for the simulators.}
    \begin{tabular}{l|lllll}
    \textbf{\begin{tabular}[c]{@{}l@{}}Configuration\\ variable\end{tabular}} & \textit{\textbf{CloudSim}} & \textit{\textbf{DISSECT-CF}} & \textit{\textbf{NetworkCloudSim}} & \textit{\textbf{OpenDC}} & \textit{\textbf{SimGrid}} \\ \hline
    Workload definition & \begin{tabular}[c]{@{}l@{}}code\end{tabular} & swf trace & code with communication & parquet trace & \begin{tabular}[c]{@{}l@{}} code\\ code with communication\end{tabular} \\ \arrayrulecolor{lightgray}\hline
    Power models & \begin{tabular}[c]{@{}l@{}}linear\\ square root\\ square\\ cubic\end{tabular} & \begin{tabular}[c]{@{}l@{}}constant\\ linear\end{tabular} & \begin{tabular}[c]{@{}l@{}}linear\\ square root\\ square\\ cubic\end{tabular} & \begin{tabular}[c]{@{}l@{}}constant\\ linear\\ square root\\ square\\ cubic\end{tabular} & linear \\ \hline
    Total configurations & 4 & 2 & 4 & 5 & 1
    \end{tabular}
    \label{tab:config-vars}
\end{table*}


The microservice applications (see \autoref{sec:apps}) and their treatments are the other independent variables. The treatments are based on previous work that involved measuring the energy consumption of \textit{Sock Shop} and \textit{Train Ticket} \cite{rodriguez_sanchez_multivariate_2024}.
They consist of two scenarios and two different user loads. The user loads are 100 and 1000, which are achieved by gradually increasing the number of users at a rate of 5 users per second.

\subsubsection{Dependent Variables}
The aim of this study is to evaluate the selected simulators for their use to simulate the \rone{host-level} energy consumption of microservice applications. For RQ1, the accuracy of the simulators and their predefined power models in terms of simulating \rone{host-level} energy consumption is measured by calculating the mean relative error (MRE) over the $N$ executions of each simulator with respect to the measured energy consumption (see \autoref{eq:relError}). The main dependent variable for this study is therefore the simulated \rone{host-level} energy consumption in Joules.

\begin{equation} \label{eq:relError}
    MRE = \frac{1}{N} \sum_i^N \frac{|\text{measured energy}_i - \text{simulated energy}_i|}{\text{measured energy}_i}
\end{equation}

For RQ2.1 concerning the customizability of the simulators, the elicited information are the properties that \textit{can} and \textit{have to} be set, for both the hardware system and workload to be simulated, as well as through what method. For RQ2.2 regarding the usability of the simulators, the mean runtime of the simulators, as well as the quality of the simulator's documentation are taken into account. This quality is assessed based on the availability of documentation, whether it is up to date, and whether it covers the energy consumption simulation features of the simulators.

\subsection{Experiment Design}
To obtain the ground truth against which the simulators are compared, the systems under test are measured with the treatments described in \autoref{sec:independentVars} using a cross-over design \cite{vegas_crossover_2016}. All possible combinations of 2 applications $\times$ 2 scenarios $\times$ 2 user loads result in 8 different treatment combinations that are executed in a randomized fashion. This is to ensure that the order of execution does not introduce bias while possible influences from the environment, network or system are not localized to one treatment. Furthermore, each treatment is executed 30 times with a 3-minute warm-up period before 5 minutes of measuring and finally a 2-minute cool down period.
